\begin{longtable}[c]{|m{0.46\textwidth}|m{0.46\textwidth}|}
    \caption{Historia de Usuario QTE-01}
    \label{tab:qte01} \\
    
    \hline
    \endfirsthead
    \endhead
    \multicolumn{2}{r}{\textit{Continúa en la siguiente página...}} \\
    \endfoot
    \endlastfoot

    \textbf{ID:} QTE-01 & \textbf{Usuario:} Desarrollador \\
    \hline
    \multicolumn{2}{|l|}{\textbf{Título:} Traducción de operaciones DML básicas} \\
    \hline
    \textbf{Prioridad en negocio:} Alta & \textbf{Riesgo en desarrollo:} Alto \\
    \hline
    \textbf{Estimación:} 8 & \textbf{Iteración asignada:} Sprint 1 \\
    \hline
    \multicolumn{2}{|p{0.92\textwidth}|}{
        \textbf{Historia de Usuario:}
        \begin{itemize}[nosep, leftmargin=0.35cm, label={}, after=\vspace{0.5em}]
            \item \textbf{Como} desarrollador,
            \item \textbf{Quiero} traducir sentencias SELECT a lenguaje Cypher,
            \item \textbf{Para} obtener datos específicos de una base de datos orientada a grafos.
        \end{itemize}} \\
    \hline
    \multicolumn{2}{|p{0.92\textwidth}|}{
        \textbf{Criterios de aceptación:} 
        \begin{itemize}[nosep, leftmargin=0.75cm, after=\vspace{0.5em}]
            \item El sistema deberá traducir  la sentencia \texttt{SELECT} para todas las columnas o solo columnas específicas de una tabla.
            \item El sistema deberá soportar la cláusula \texttt{WHERE}, operadores de comparación (\texttt{=}, \texttt{<}, \texttt{<=}, \texttt{>}, \texttt{>=}, \texttt{<>}) y operadores lógicos (\texttt{AND}, \texttt{OR} y \texttt{NOT}).
            \item El sistema deberá traducir la cláusula \texttt{ORDER BY} junto a la forma de ordenamiento (\texttt{ASC} o \texttt{DESC}).
            \item El sistema deberá traducir correctamente las palabras reservadas \texttt{DISTINCT}, y \texttt{LIMIT/TOP}.
            \item El usuario deberá ingresar una sentencia SQL que cumpla con las restricciones indicadas.
            \item El usuario deberá seleccionar la opción \textbf{Traducir} para visualizar la sentencia en lenguaje Cypher equivalente a la que ingresó.
            \item Si la sentencia tiene errores de sintaxis o no es soportada, el sistema indicará un mensaje de error y no podrá ser traducida.
            \item Las consultas traducidas, serán agregadas automáticamente al historial.
        \end{itemize}} \\
    \hline
\end{longtable}